% Рисунок с доски (двумерное пространство-время): треугольник зависимости точки
% (x,t), основание -- интервал B_{x,at} = (x-at, x+at).
\begin{figure}[htbp]
\centering
\begin{tikzpicture}[>={Latex[length=4pt]}]
  \draw[->] (-0.4,0) -- (6.0,0) node[anchor=north east] {$y$};
  \draw[->] (0,-0.4) -- (0,3.4) node[anchor=south east] {$t$};

  \fill[pattern=north east lines, pattern color=blue!35] (1.0,0) -- (5.0,0) -- (3.0,2.6) -- cycle;
  \draw[thick, blue!60!black] (1.0,0) -- (5.0,0) -- (3.0,2.6) -- cycle;

  \fill (3.0,2.6) circle (1.6pt);
  \node[anchor=south, font=\footnotesize] at (3.0,2.68) {$(x,t)$};
  \node[anchor=west, font=\footnotesize] at (4.1,1.6) {$|y - x| = a(t-\tau)$};

  \draw[dashed, gray] (0,1.3) -- (2.0,1.3);
  \draw[very thick, red!60!black] (2.0,1.3) -- (4.0,1.3);
  \node[anchor=east, font=\footnotesize] at (-0.08,1.3) {$\tau$};
  \node[red!60!black, anchor=south, font=\footnotesize] at (3.0,1.34) {$\Bxat{x}{a\tau'}$};

  \draw[very thick, blue!60!black] (1.0,0) -- (5.0,0);
  \node[anchor=north, font=\footnotesize] at (1.0,-0.05) {$x-at$};
  \node[anchor=north, font=\footnotesize] at (5.0,-0.05) {$x+at$};
  \node[anchor=north, font=\footnotesize] at (3.0,-0.05) {$x$};
  \node[blue!60!black, anchor=north, font=\footnotesize] at (3.0,-0.55)
    {$\Bxat{x}{at}$ — область зависимости};
\end{tikzpicture}
\caption{Двумерное пространство-время $(y,t)$. Значение $u(x,t)$ определяется
данными только внутри заштрихованного характеристического треугольника; его
основание — интервал $\Bxat{x}{at} = (x-at,\,x+at)$. Здесь
$\tau' = t - \tau$ — «оставшееся» время, так что срез на уровне $\tau$ есть
$\Bxat{x}{a(t-\tau)}$.}
\label{fig:dependence}
\end{figure}
